\section{The Books Were Opened}\label{sec:books}

The Apocalypse gives the last judgment its most enduring image:

\begin{quote}
And I saw a great white throne and one sitting upon it, from whose face the earth and heaven fled away: and there was no place found for them. And I saw the dead, great and small, standing in the presence of the throne. And the books were opened: and another book was opened, which is the book of life. And the dead were judged by those things which were written in the books, according to their works. (Apoc 20:11--12)
\end{quote}

The \latin{Dies irae} takes up the image in two strophes:

\begin{quote}
\latin{Liber scriptus proferetur,\\
\quad In quo totum continetur,\\
\quad Unde mundus judicetur.\\
Judex ergo cum sedebit,\\
\quad Quidquid latet, apparebit:\\
\quad Nil inultum remanebit.}

\medskip
Lo! the Book exactly worded,\\
Wherein all hath been recorded;\\
Thence shall judgment be awarded.\\
When the Judge His seat attaineth,\\
And each hidden deed arraigneth,\\
Nothing unavenged remaineth.\endnote{Latin: \emph{Missale Romanum} (Desclée, 1889), p.~92*. English: Irons, stanza~3, \emph{Hymns Ancient and Modern}, hymn 398, p.~307.}
\end{quote}

What is this book in which everything is written? Augustine, reading the chapter closely in the twentieth book of \emph{The City of God}, understood the plural books to be the Scriptures of the Old and New Testaments, ``that out of them it might be shown what commandments God had enjoined,'' and the single book, the book of each man's life, to show ``what commandments each man has done or omitted to do.'' Then he asks how such a book could be read:

\begin{quote}
If this book be materially considered, who can reckon its size or length, or the time it would take to read a book in which the whole life of every man is recorded? Shall there be present as many angels as men, and shall each man hear his life recited by the angel assigned to him? \ldots\ We must therefore understand it of a certain divine power, by which it shall be brought about that every one shall recall to memory all his own works, whether good or evil, and shall mentally survey them with a marvellous rapidity, so that this knowledge will either accuse or excuse conscience, and thus all and each shall be simultaneously judged. And this divine power is called a book, because in it we shall as it were read all that it causes us to remember.\endnote{Augustine, \emph{De civitate Dei} XX.14 (Dods).}
\end{quote}

The book, then, is memory, opened by God. Nothing will be read to the soul from outside that it does not already carry within it. Every word, every act, every omission, every secret consent, everything that the passing of years has buried and the habit of excuse has softened, will be recalled at once and seen as it is. Paul had said as much of the Gentiles, ``their conscience bearing witness to them: and their thoughts between themselves accusing or also defending one another, in the day when God shall judge the secrets of men by Jesus Christ'' (Rom 2:15--16). Christ himself extends the account to the smallest coin of human life: ``every idle word that men shall speak, they shall render an account for it in the day of judgment. For by thy words thou shalt be justified, and by thy words thou shalt be condemned'' (Matt 12:36--37). And he promises that the account will be public: ``For there is nothing covered that shall not be revealed: nor hidden that shall not be known. For whatsoever things you have spoken in darkness shall be published in the light'' (Luke 12:2--3).

Alphonsus, following the ancient reading, speaks of two books at the particular judgment: ``There shall be two books, the Gospel and conscience. In the Gospel shall be read what the accused should have done, and in his conscience, what he has done.''\endnote{Alphonsus Liguori, \emph{Preparation for Death}, Consideration~XXIV, second point (Duffy, 1843), p.~293.} The two together make the whole judgment: the law of Christ, and the life measured against it.

\subsection{Reading the book now}

The saints drew from this doctrine a practice that the Church has never ceased to commend: the examination of conscience. If the book of one's life will be opened at death, it can be opened beforehand, by the soul itself, in the presence of God, while there is still time to write in it a different ending. Francis de Sales, in the meditation on judgment that he set in the first part of the \emph{Introduction to the Devout Life}, leads the reader through the separation of the good from the wicked and then says:

\begin{quote}
This separation made, the books of conscience will be opened, and all men will behold the malice of the wicked, and how they have contemned God; as also the penitence of the good, and the results of the grace they received. Nothing will be hid. O my God, what confusion to the one, what rejoicing to the other!\endnote{Francis de Sales, \emph{Introduction to the Devout Life}, part~I, chapter~14, ``Sixth Meditation: On Judgment,'' consideration~5 (London: Rivingtons, 1876), p.~39.}
\end{quote}

Notice what the good have to show: not their innocence, but ``the penitence of the good, and the results of the grace they received.'' The book of the saved is not a blank page. It is a page on which sins are written and also their repentance, and over both the grace of God that brought the one out of the other. De Sales draws the conclusion in the affections that follow:

\begin{quote}
Abhor your sins, which alone can cause you to be lost when that fearful day comes. Surely I will judge myself now, that I be not judged;---I will examine my conscience, accuse, condemn, punish myself, that the Judge may not condemn me then. I will confess my faults, and follow the counsels given me.\endnote{De Sales, \emph{Introduction}, part~I, chapter~14, Affections and Resolutions, 2 (Rivingtons, 1876), p.~40.}
\end{quote}

His instruction is to end the meditation by thanking God ``for having given you means of safety in that terrible Day, and time for repentance.''\endnote{De Sales, \emph{Introduction}, part~I, chapter~14, Conclusion (Rivingtons, 1876), p.~40.} The time for repentance is the present. The book is still open, and the hand that writes in it is still one's own.

Augustine's image is also a comfort. A book written by an accuser might contain lies; the book God opens at the judgment contains only the truth, and the truth includes every act of contrition, every confession, every sin that the blood of Christ has washed. That is why de Sales can say that the good will be seen in their penitence and in the fruits of grace: what has been confessed and forgiven stands in the book as forgiven. The soul that has learned to read its own book honestly before the Lord, day by day, will find at the judgment that nothing in it is a surprise, and that the last line on every page is mercy.
